\section{Formalization}\label{sec:lean}

The definitions D1--D6, the general statements of E1--E16 and the laws and instances of
\Cref{sec:concrete} are checked in Lean~4 \citep{DeMouraUllrich2021}, in a library with no external
dependencies beyond the vendored interfaces of the earlier papers of the series. The library builds
without \texttt{sorry} and declares no axioms. Following \citet{Pollack1998}, what a reader must
trust is the kernel and the statements; \Cref{supp:concordance} maps every numbered result of this
paper to its Lean declarations, file and line, so that each statement can be compared with the
printed one.

\paragraph{Axioms.} Every declaration uses at most the three standard axioms \texttt{propext},
\texttt{Quot.sound} and \texttt{Classical.choice}; the record for each theorem is in
\Cref{supp:status}.

\paragraph{What each checked statement covers.}
\begin{itemize}
  \item Definitions: \Cref{prop:d1,prop:d2} and \Cref{lem:d3,lem:d4,lem:d5,lem:d6}. In
    \Cref{prop:d2} the chain rule and positive semidefiniteness are hypotheses.
  \item Laws and instances: \Cref{prop:device-step,thm:e3-episode,prop:e2-self,thm:e2-declared,thm:e2-tower,prop:e1-forcing,thm:e1-clocked,prop:e12-structure,thm:e12-repair,prop:e12-overlap,thm:e13-chain,thm:e13-cross,prop:e13-flat}.
    The converse of \Cref{thm:e3-episode} is checked for the form of the definition that requires
    only the move entry. \Cref{prop:e2-self} is checked in the weaker form that the classifier does
    not accept the claim. In \Cref{thm:e13-chain} the payment of the composite is the sum by
    definition.
  \item General theorems: \Cref{thm:e6-kkt,thm:e6-caps,thm:e9,prop:e9-saturation,prop:e15-identity,prop:e4-brittle},
    with the numerical instance \Cref{thm:e6-atcap} of \Cref{thm:e6-caps}. The
    identity of \Cref{prop:e15-identity} is derived from the two exhaustion equations and the gate.
  \item Certificate specifications and definitions: the Templates of \Cref{sec:templates} and
    \Cref{supp:templates},
    \Cref{prop:e12-singleton}, and \Cref{lem:e5,lem:e8}, each exactly as printed: the certificate is
    a hypothesis, and the conclusion is derived from it.
\end{itemize}

\paragraph{Argued on paper only.} \Cref{prop:e6-sufficient,ex:e6-nonconcave,rem:e6-unlinked}; the
envelope reading of E8; alarm fatigue (E7.1); and operational selfhood (E12.1).

\paragraph{Declared propositions.} Formedness of a package and completeness of record, ledger and
history inventories are propositions supplied with the data. Lean checks what follows from them,
not that they hold of a real system. In the repair device and the clocked device they are declared
true, which is why their instances are weak. The instance of E6 is numerical and carries no ledger
records.
